\section{Discussion}\label{sec:discussion}

\paragraph{Limitations.}
The fleet writes code but never pushes: integration to shared remotes
remains human-gated, so measured autonomy covers worktree-to-acceptance,
not publication. Gate scripts themselves are authored by agents and
could be gamed; the protocol mitigates this by pre-registering criteria
append-only before execution and requiring the verifier role to run
gates from a different mission than the implementer that wrote the
code under test. The relay models are a convenience sample of what a
production relay exposes; findings are about the \emph{protocol}, not a
leaderboard of models.

\paragraph{Safety posture.}
Three properties hold by construction: the agent platform account is
read-only (dispatch authority derives from claim tokens, not role
elevation); no worktree can push to a remote; and every promotion and
acceptance decision terminates with a human owner. Envelope parameters
are policy, tunable by the owner without code changes.

\paragraph{Reproducibility.}
Mission rows, event streams, approval chains, run manifests, and
artifact checksums are exported directly from the platform database;
the paper's empirical sections will cite only exported artifacts,
never hand-copied numbers. Campaign artifacts and analysis scripts
accompany the submission.

\paragraph{Research-process framing.}
Treating development as research is not rhetorical: each mission is a
unit with pre-registered criteria, bounded budget, immutable event
trail, and adjudicated outcome. The same substrate that audits the
physics audits the fleet, which is precisely what makes RQ2 answerable
without new instrumentation.
